\documentclass[aspectratio=169]{beamer}
\usepackage{will_handley_beamer}
\usepackage{title_page}

% Commands
% --------
% - \arxiv{arxiv number}
% - \arxiv{<number>}            arxiv.org/abs/<number>
% - \oldarxiv{<arxiv number>}   arxiv.org/<number>
% - \doi{<doi>}                 doi.org/<doi>
% - \xkcd{<number>}             xkcd.com/<number>
% - \email{<email>}             <<email>>
% - \tthref{<website>}          <website>
% - \av[dist]{<quantity>}       <quantity>_{dist}
% - \student{<name>}{<detail>}{<photo>}

% Talk details
% ------------
\title{AI/ML Tools for Research}
\subtitle{KICC Workshop}
\date{23rd July 2025}

\begin{document}

\begin{frame}
    \titlepage
\end{frame}

\begin{frame}
    \frametitle{Overview}
    
    \begin{columns}
        \column{0.5\textwidth}
        \begin{itemize}
            \item Follow-up to successful Part II/III training.
            \item Focus on ``doing'' not ``explaining''.
            \item Three layers of AI tools framework.
            \item Interactive demonstrations.
            \item 90 minutes then tea.
        \end{itemize}
        
        \column{0.5\textwidth}
        \textbf{Common questions:}
        \begin{itemize}
            \item ``Which model should I use?''
            \item ``Don't these things hallucinate?''
            \item ``How do I get started?''
            \item ``Is this worth the investment?''
        \end{itemize}
    \end{columns}
\end{frame}

\begin{frame}
    \frametitle{First things first}
    If you've been busy, the past six months have seen a couple of inflection points
    \vspace{0.5cm}
    \begin{itemize}
        \item \textbf{February 2025}: Models acquired many of the capabilities we prize in our PhD students
            \begin{itemize}
                \item Models: o3 series, Gemini 2.5, Claude 3.5 Sonnet.
                \item Capabilities: Code development, mathematical reasoning, literature review, paper drafting, grant writing.
            \end{itemize}
        \item \textbf{May 2025}: Agentic systems launched commercially
            \begin{itemize}
                \item Tools: Claude Code, Cursor agent mode, Deep Research.
                \item Capabilities: Writing test suites, assembling pip-installable codes, synthesising handwritten notes and TeX files, write/run/debug code in languages you don't know.
            \end{itemize}
        \item \textbf{The performance gap}: Free models lag significantly behind paid versions
    \end{itemize}
    
    \vspace{0.5cm}
    If your lived experience is typing a research question into free ChatGPT at the start of the year and finding it was rubbish, things have moved on
\end{frame}

\begin{frame}
    \frametitle{On AI Hype}
    \begin{columns}
        \column{0.48\textwidth}
        \begin{itemize}
            \item These things are not ``intelligent'' in the way Silicon Valley wants you to believe (AGI is not just around the corner).
            \item \emph{Caveat emptor}: It is very hard for a human to read language, and not reflexively construct a mind behind it.
            \item You should not however think of them as mere ``next token predictors'':
                \begin{itemize}
                    \item The modern reinforcement learnt reasoning systems are significantly more than that.
                    \item Agentic tools ground these in reality.
                \end{itemize}
        \end{itemize}
        
        \column{0.48\textwidth}
        \begin{block}{A senior examiner's experiment}
            \begin{itemize}
                \item I got o3 to take the Part II exams.
                \item Prompt on exam day: ``You are a first class Part II Cambridge astronomy student... here is the syllabus... here is the question, answer it''.
                \item Toby Lovick transcribed answers onto written exam scripts.
                \item Slipped them amongst real scripts and marked blind.
                \item \textbf{Result}: Best student we've had across IoA history, even without coursework.
            \end{itemize}
        \end{block}
    \end{columns}
\end{frame}

\begin{frame}
    \frametitle{The AI landscape}
    \centering
    \begin{tabular}{l|c|c|c|}
        \visible<3->{Under the hood} & 
        \visible<3->{\includegraphics[valign=m, width=0.2\textwidth]{logos/microsoft}} & 
        \visible<3->{\includegraphics[valign=m, width=0.2\textwidth]{logos/amazon}} & 
        \visible<3->{\includegraphics[valign=m, width=0.2\textwidth]{logos/google}} \\[0.5cm]
        
        \visible<2->{Companies} & 
        \visible<2->{\includegraphics[valign=m, width=0.2\textwidth]{logos/openai}} & 
        \visible<2->{\includegraphics[valign=m, width=0.2\textwidth]{logos/anthropic}} & 
        \visible<2->{\includegraphics[valign=m, width=0.2\textwidth]{logos/deepmind}} \\[0.5cm]
        
        \visible<1->{Models} & 
        \visible<1->{\includegraphics[valign=m, width=0.2\textwidth]{figures/chatgpt.png}} & 
        \visible<1->{\includegraphics[valign=m, width=0.2\textwidth]{figures/claude.pdf}} & 
        \visible<1->{\includegraphics[valign=m, width=0.2\textwidth]{figures/gemini.pdf}}
    \end{tabular}
    
    \vspace{1cm}
    \visible<4->{
        \textbf{Others exist:} Perplexity, Poe, Character.ai, You.com, Cohere, xAI, Llama...
    }
    
    \vspace{0.5cm}
    \visible<5->{
        \textbf{But these are the three main pillars (July 2025).}
    }
\end{frame}

\begin{frame}
    \frametitle{Three layers of AI tools}
    \framesubtitle{A framework for understanding the landscape}
    \begin{block}{Layer 3: Agentic Systems}
        \vspace{-0.3cm}
        \begin{columns}
            \column{0.48\textwidth}
            \begin{itemize}
                \item Claude Code, cursor agent mode.
                \item Custom workflows.
            \end{itemize}
            \column{0.48\textwidth}
            \begin{itemize}
                \item Autonomous task completion.
                \item Steeper learning curve but transformative.
            \end{itemize}
        \end{columns}
    \end{block}
    
    \vspace{0.5cm}
    
    \begin{block}{Layer 2: Prompt Engineering}
        \vspace{-0.3cm}
        \begin{columns}
            \column{0.48\textwidth}
            \begin{itemize}
                \item ChatGPT, Claude, Gemini.
                \item Web interfaces.
            \end{itemize}
            \column{0.48\textwidth}
            \begin{itemize}
                \item Interactive problem solving.
                \item Good for exploration and learning.
            \end{itemize}
        \end{columns}
    \end{block}
    
    \vspace{0.5cm}
    
    \begin{block}{Layer 1: Autocomplete}
        \vspace{-0.3cm}
        \begin{columns}
            \column{0.48\textwidth}
            \begin{itemize}
                \item GitHub Copilot.
                \item VS Code extensions.
            \end{itemize}
            \column{0.48\textwidth}
            \begin{itemize}
                \item Completes code as you type.
                \item Minimal learning curve.
            \end{itemize}
        \end{columns}
    \end{block}
\end{frame}

\begin{frame}
    \frametitle{Layer 1: Autocomplete}
    \framesubtitle{Standard code-completion, but powered by AI}
    
    \begin{columns}
        \column{0.6\textwidth}
        \begin{block}{Core idea}
            \begin{itemize}
                \item AI-powered code-completion.
                \item Trained on all of GitHub.
                \item Context-aware suggestions.
            \end{itemize}
        \end{block}
        
        \begin{block}{Recommendation}
            \begin{itemize}
                \item GitHub Copilot (Pro subscription).
                \item Free for university email holders.
                \item \url{github.com/settings/education/benefits}.
                \item Worth \$10/month.
                \item Also available in Cursor.
            \end{itemize}
        \end{block}

        \column{0.4\textwidth}
        \includegraphics[width=\textwidth]{figures/copilot_screenshot.png}
    \end{columns}
\end{frame}

\begin{frame}
    \frametitle{Layer 2: Prompt Engineering}
    \framesubtitle{The usual chat versions}
    
    \begin{columns}
        \column{0.3\textwidth}
        \centering
        \includegraphics[width=\textwidth]{figures/chatgpt_screenshot.png}
        
        ChatGPT
        
        \url{chat.openai.com}
        
        \column{0.3\textwidth}
        \centering
        \includegraphics[width=\textwidth]{figures/claude_screenshot.png}
        
        Claude
        
        \url{claude.ai}
        
        \column{0.3\textwidth}
        \centering
        \includegraphics[width=\textwidth]{figures/gemini_screenshot.png}
        
        Gemini
        
        \url{gemini.google.com}
    \end{columns}
\end{frame}

\begin{frame}
    \frametitle{Layer 2: Developer versions}
    \framesubtitle{Use their developer level tools}
    
    \begin{columns}
        \column{0.3\textwidth}
        \centering
        \includegraphics[width=\textwidth]{figures/chatgpt_dev_screenshot.png}
        
        OpenAI Playground
        
        \url{platform.openai.com/playground}
        
        \column{0.3\textwidth}
        \centering
        \includegraphics[width=\textwidth]{figures/claude_dev_screenshot.png}
        
        Claude Workbench
        
        \url{console.anthropic.com/workbench}
        
        \column{0.3\textwidth}
        \centering
        \includegraphics[width=\textwidth]{figures/gemini_dev_screenshot.png}
        
        Google AI Studio
        
        \url{aistudio.google.com}
    \end{columns}
    
    \vspace{0.3cm}
    
    \textbf{These should be your default as a scientist} -- if you're happy using a command line, you shouldn't be satisfied by an interface that looks like a child's toy.
\end{frame}

\begin{frame}
    \frametitle{Layer 2: Prompt Engineering}
    \framesubtitle{Use AI to improve AI}
    \begin{columns}
        \column{0.48\textwidth}
        \begin{itemize}
            \item Likely the layer of AI most familiar to you.
            \item ``chatting''/conversation can be very powerful for naturally tuning the attention/context of the model.
            \item For one shot work, a little prompt engineering goes a long way:
                \begin{itemize}
                    \item You can use these tools to improve prompts.
                \end{itemize}
        \end{itemize}
        \column{0.48\textwidth}
        \includegraphics<1>[width=\textwidth]{figures/prompt_demo1.png}%
        \includegraphics<2>[width=\textwidth]{figures/prompt_demo2.png}
    \end{columns}
\end{frame}

\begin{frame}
    \frametitle{Layer 2: Cost}
    \framesubtitle{Pay as you use vs subscription models}
    
    \begin{columns}
        \column{0.6\textwidth}
        \begin{itemize}
            \item Gemini is free to use (but you sign away your data to Google).
            \item ChatGPT and Claude offer \$20/month subscriptions.
            \item API access available: pay per token consumed:
                \begin{itemize}
                    \item \$0.0002-0.015 per 1000 tokens (roughly 750 words).
                    \item Orders of magnitude cheaper for research use.
                    \item Only pay for what you actually use.
                \end{itemize}
            \item Never buy yearly subscriptions - AI evolves too rapidly.
            \item My group's preference: pay for API access via tokens.
                \begin{itemize}
                    \item Pay for what you use.
                    \item Allows you to try out the latest models.
                    \item API keys allow you to do all kinds of tricks at the command line.
                \end{itemize}
        \end{itemize}
        
        \column{0.4\textwidth}
        \includegraphics[width=\textwidth]{figures/cost_demo1.png}
        \includegraphics[width=\textwidth]{figures/cost_demo2.png}
    \end{columns}
\end{frame}

\begin{frame}
    \frametitle{Layer 3: Agentic Systems}
    \framesubtitle{Autonomous task completion - transformative but complex}
    
    \begin{columns}
        \column{0.6\textwidth}
        \begin{itemize}
            \item Emerged commercially in April 2025.
            \item If you get too skilled at prompt engineering, your job ends up being a copy-paste-run-debug monkey.
            \item Agentic systems are what happens if you let ChatGPT edit files, run commands, and inspect the behaviour.
            \item This closes the scientific loop.
            \item Agentic systems are \textbf{grounded in reality}. They self-correct hallucinations by testing their hypotheses.
            \item Claude Code (CLI), Cursor Agent mode, Gemini CLI.
            \item Deep research also falls into this category (around for longer)
            \item Goal-based vs task-based.
        \end{itemize}
        
        \column{0.4\textwidth}
        \includegraphics<1>[width=\textwidth]{figures/agentic_demo.png}%
        \includegraphics<2>[width=\textwidth]{figures/agentic_demo2.png}%
    \end{columns}
\end{frame}

\begin{frame}
    \frametitle{Layer 3: Agentic Systems Cost}
    \framesubtitle{More expensive but potentially more valuable}
    
    \begin{itemize}
        \item Significantly more expensive than prompt engineering.
        \item Can consume hundreds of API calls per task.
        \item For agentic systems, subscriptions \emph{are} cheaper than pay-per-token.
        \item I was spending \$40/day on Claude Code with API tokens.
        \item Subscription gives you about 2hr coding session per day for \$20/month/user.
        \item Gemini CLI is ``free'' but risks going over free tier limit, and is less stable.
        \item In July 2025, Claude Code best in class.
        \item There are many further unknown costs of shifting your research toward ``vibe coding''.
    \end{itemize}
\end{frame}

\begin{frame}
    \frametitle{Conclusions and Getting Started}
    
    \begin{columns}
        \column{0.5\textwidth}
        \begin{itemize}
            \item Use developer-level tools, not consumer interfaces.
            \item API access typically cheaper for research use.
            \item Capturing and owning content is critical -- use Otter for transcription.
            \item Export conversations and notes in formats that can be fed to LLMs.
            \item Get into the habit of recording information in exportable forms.
            \item Agentic systems are transformative but carry additional risks requiring careful consideration.
        \end{itemize}
        
        \column{0.5\textwidth}
        \begin{block}{To-do}
            \begin{itemize}
                \item Set up GitHub Copilot Pro (takes a couple of days to approve).
                \item Try AI Studio rather than ChatGPT.
                \item Get your PI to set up some tokens for your group to use (you only need ~O(\$10) to try it out).
                \item Try Google Gemini's Deep Research \textbf{button}.
                \item Test Claude Code/Gemini CLI for a month.
            \end{itemize}
        \end{block}
    \end{columns}
\end{frame}

\end{document}
